\section{The Optimistic Case}
\label{sec:optimistic}

In this section, we show that choosing an optimistic online learner \citep{chiang2012online, rakhlin2013online, steinhardt2014eg} will accelerate our DP online-to-batch algorithm. Optimistic algorithms  are provided with additional ``hints'' in the form of $\hat \ell_t(w)=\langle \hat g_t, w\rangle$ as an approximation of the true loss $\ell_t(w)=\langle g_t,w\rangle$, and they can incorporate $\hat \ell_t$ to decide $w_t$. The regret of an optimistic algorithm depends on the quality of hints: if $\hat{g}_t\approx g_t$, then it achieves a low regret.  Formally, in this paper, we say an online learning algorithm is \textit{optimistic w.r.t. norm $\|\cdot\|$} if its regret is the following:
\begin{equation}
    \regret_T(x^*) \leq O\left( D\sqrt{\sum_{t=1}^T \| \hat{g}_t - g_t\|_*^2 } \right), \label{eq:optimistic-OL}
\end{equation}
where $D$ denotes the diameter of the learner's domain.

A common choice of the hint $\hat{g}_t$ is $g_{t-1}$, the gradient in the last round since intuitively, one could expect $g_{t-1}\approx g_t$ when the loss functions are smooth. In this section, we also follow this choice. Recall that in Algorithm \ref{alg:DP-OTB}, the online learner receives $t$-th loss $\ell_t(w) = \langle g_t+\gamma_t, w\rangle$, where $g_t=\sum_{i=1}^t\delta_i$ is the sum of gradient differences, and $\gamma_t$ is some noise. Therefore, we define the $t$-th hint as $\hat{g}_t = g_{t-1}+\gamma_{t-1}$.


\begin{theorem}
\label{thm:DP-OTB:optimistic}
% Suppose $\|\cdot\|^2$ is $\lambda$-strongly convex, $W$ is bounded by $D$, and $\ell$ is $G$-Lipschitz and $H$-smooth. 
Suppose Assumption \ref{as:DP-OTB:sc-norm} - \ref{as:DP-OTB:smooth} hold, and $\calD$ is a $(V,\alpha)$-RDP distribution. Set $\weight_t=t^k$ and $\sigma_t^2$ as defined in \eqref{eq:DP-OTB:RDP-noise},
If the online learner is optimistic (satisfying \eqref{eq:DP-OTB:optimistic-OL}) with $t$-th gradient $\bar{g}_t=g_t+\gamma_t$ and $t$-th hint $\hat{g}_t=\bar{g}_{t-1}$, then
\begin{align*}
    \frac{\E[\regret_T(x^*)]}{\weight_{1:T}} 
    % \leq \sqrt{3}(k+1)^2D(G+DH)\left(1+\frac{4\sqrt{V}\log_2(2T)}{\sqrt{\lambda}\rho}\right)\frac{1}{T^{3/2}}.
    \leq O\left( D(G+DH)\left(1+\frac{\sqrt{V}\log T}{\sqrt{\lambda}\rho}\right)\frac{1}{T^{3/2}} \right).
\end{align*}
% Consequently, if we further assume Assumption \ref{as:sigma-G} and \ref{as:sigma-H}, then
% \begin{align*}
%     \E[f(x_T) - f(x^*)]
%     &\leq \frac{2(k+1)^2D}{\sqrt{\lambda}}\left(\frac{\sigma_G+D\sigma_H}{\sqrt{T}}+\frac{2\sqrt{K}(G+DH)\log_2^{3/2}T}{\rho T}\right).
% \end{align*}
\end{theorem}

The proof is in Appendix \ref{app:optimistic}.
As an immediate corollary, if we further assume Assumption \ref{as:sigma-G} and \ref{as:sigma-H}, then Theorem \ref{thm:DP-OTB-general-norm} applies. Together with this theorem, they imply that optimistic learners achieve the following convergence rate that is adaptive to the variance:
\begin{align*}
    \E[f(x_T)-f(x^*)] = O\left( \frac{D(\sigma_G+D\sigma_H)}{\sqrt{T}} + \frac{\sqrt{d}D(G+DH)\log T}{\rho T} \right).
\end{align*}
Compared to the bound in the non-optimistic case (Remark \ref{rmk:base-result}), this bound has $\sigma_G$ instead of $G$ in the first term. Thus, when the gradient $\nabla\ell(x_t,z_t)$ has low variance, i.e., $\sigma_G\ll G$, the optimistic bound outperforms the standard bound.